\chapter{Доверителни интервали и~хипотези}

\section{Доверителни интервали}

\begin{rembox}[Постановка]
Нека $X$ е от клас $F_X(x;\theta)$, $\theta=(\theta_1,\dots,\theta_s)$,
и~$\bar X=(X_1,\dots,X_n)$ е извадка. Търсим $I_1=I_1(\bar X) < I_2
(\bar X)=I_2$, такива че
\[
  \mathbb P(I_1\le\theta\le I_2) \ge \gamma,
\]
където $\gamma\in(0,1)$ (обикновено $\gamma=0.9,0.95,0.99$) се нарича
\textbf{ниво на доверие} и~зависи от най-голямата вероятност за
грешка, която можем да си позволим да допуснем. \emph{Важно:} $I_1$
и~$I_2$ са стохастични (зависят от случайната извадка), докато
$\theta$ е фиксирано - търсим интервал (със случайни краища), който
да „улавя“ фиксирания параметър с~висока вероятност.
\end{rembox}

\begin{defbox}[~(Централна статистика)]
Казваме, че $T=T(\bar X,\theta)$ е \textbf{централна статистика},
ако:
\begin{enumerate}[label=\arabic*)]
  \item $T(\bar X,\cdot)$ е монотонна функция на $\theta$;
  \item разпределението на $T(\bar X,\theta)$ не зависи от $\theta$, т.е.\ $\mathbb P(T(\bar X,\theta)<x)=G(x)$ за някаква функция $G$, зависеща само от $x$.
\end{enumerate}
(Централните статистики зависят от $\theta$ като аргумент, но
разпределението им - не.)
\end{defbox}

\begin{rembox}[Конструкция на доверителен интервал от централна статистика]
Нека $T$ е монотонно растяща по $\theta$ (случаят намаляваща е
аналогичен, само с~разменени краища). Избираме квантили $q_1<q_2$
такива, че $G(q_2)-G(q_1)=\gamma$ (например симетрично, ако $G$ е
симетрична относно нулата). Тогава
\[
  \gamma = \mathbb P\big(q_1\le T(\bar X,\theta)\le q_2\big) = G(q_2)-G(q_1).
\]
Понеже $T(\bar X,\cdot)$ е растяща (следователно обратима по
$\theta$), можем да „обърнем“ неравенствата, прилагайки $T^{-1}(
\bar X,\cdot)$ към трите части:
\[
  \gamma = \mathbb P\big(T^{-1}(\bar X,q_1) \le \theta \le T^{-1}(\bar X,q_2)\big).
\]
Следователно $I_1=T^{-1}(\bar X,q_1)$, $I_2=T^{-1}(\bar X,q_2)$ (или,
ако $T$ е намаляваща по $\theta$: $I_1=T^{-1}(\bar X,q_2)$,
$I_2=T^{-1}(\bar X,q_1)$).
\end{rembox}

\begin{defbox}[~(Доверителен интервал)]
Интервалът $\big(I_1(\bar X),I_2(\bar X)\big)$, конструиран по
горния начин от централна статистика $T$, се нарича
\textbf{доверителен интервал} за $\theta$ с~ниво на доверие $\gamma$.
\end{defbox}

\section{Доверителни интервали за средната стойност на нормално разпределение}

\begin{thmbox}[~(доверителен интервал за \texorpdfstring{$\mu$}{мю} при известно \texorpdfstring{$\sigma^2$}{сигма^2})]{Теорема}
Нека $X\sim N(\mu,\sigma^2)$, $\sigma^2$ известно, и~$\bar X=(X_1,
\dots,X_n)$. Нека $\hat\mu=\overline X_n^{(1)}=\tfrac1n\sum_{j=1}^n
X_j$. Тогава доверителният интервал за $\mu$ с~ниво на доверие
$\gamma$ е
\[
  \Big(\hat\mu - q\,\frac{\sigma}{\sqrt n},\ \hat\mu+q\,\frac{\sigma}{\sqrt n}\Big), \qquad q = \Phi^{-1}\Big(\frac{1+\gamma}2\Big),
\]
където $\Phi$ е функцията на разпределение на $N(0,1)$.
\end{thmbox}

\begin{proofbox}
Понеже $X_j\sim N(\mu,\sigma^2)$ независими, $\hat\mu=\tfrac1n
\sum_j X_j \sim N(\mu,\sigma^2/n)$ (Тема~11: сума на независими
нормални случайни величини), следователно
\[
  T = \frac{\hat\mu-\mu}{\sigma/\sqrt n} \sim N(0,1)
\]
е централна статистика (разпределението ѝ, $N(0,1)$, не зависи от
$\mu$). Избираме симетрични квантили $\pm q$ на $N(0,1)$ с~$\mathbb
P(-q\le Z\le q)=\gamma$, т.е.\ $\Phi(q)-\Phi(-q)=2\Phi(q)-1=\gamma
\Rightarrow q=\Phi^{-1}\big(\tfrac{1+\gamma}2\big)$ (използвахме
$\Phi(-q)=1-\Phi(q)$ по симетрията на $N(0,1)$). Тогава
\[
  \gamma = \mathbb P\Big(-q \le \frac{\hat\mu-\mu}{\sigma/\sqrt n} \le q\Big) = \mathbb P\Big(-q\frac{\sigma}{\sqrt n} \le \hat\mu-\mu \le q\frac{\sigma}{\sqrt n}\Big) = \mathbb P\Big(\hat\mu-q\frac\sigma{\sqrt n} \le \mu \le \hat\mu+q\frac\sigma{\sqrt n}\Big). \qedbg
\]
\end{proofbox}

\begin{thmbox}[~(разлагане на квадратичната сума)]{Теорема}
Нека $X\sim N(\mu,\sigma^2)$ и~$\bar X=(X_1,\dots,X_n)$, $\hat\mu=
\overline X_n$, $S^2=\tfrac1{n-1}\sum_{j=1}^n(X_j-\hat\mu)^2$. Тогава
$S^2$ е неизместена оценка за $\sigma^2$, $S^2$ и~$\hat\mu$ са
независими, и~\[
  \frac{(n-1)S^2}{\sigma^2} \sim \chi^2(n-1).
\]
\end{thmbox}

\begin{proofbox}
Нека $Z_j=\tfrac{X_j-\mu}\sigma\sim N(0,1)$, i.i.d. По алгебричната
идентичност (разлагане чрез добавяне и~изваждане на $\hat\mu$):
\[
  \sum_{j=1}^n (X_j-\mu)^2 = \sum_{j=1}^n (X_j-\hat\mu)^2 + n(\hat\mu-\mu)^2
  \quad\Longrightarrow\quad
  \sum_{j=1}^n Z_j^2 = \frac{(n-1)S^2}{\sigma^2} + n\frac{(\hat\mu-\mu)^2}{\sigma^2}.
\]
Лявата страна $\sum_j Z_j^2\sim\chi^2(n)$ (сума на квадратите на $n$
независими $N(0,1)$-величини, Тема~11), а~вторият член отдясно е
$\big(\tfrac{\hat\mu-\mu}{\sigma/\sqrt n}\big)^2 \sim \chi^2(1)$ (по
предната теорема, изразът в~скоби е $N(0,1)$).

За да установим, че $(n-1)S^2/\sigma^2\sim\chi^2(n-1)$ и~че тя е
независима от $\hat\mu$, използваме следния класически похват
(ортогонална трансформация на нормален вектор): нека $\bar Z=(Z_1,
\dots,Z_n)$. Съществува ортогонална матрица $O$ (т.е.\ $OO^\top=I$),
чийто първи ред е $\big(\tfrac1{\sqrt n},\dots,\tfrac1{\sqrt n}\big)$
- останалите редове допълваме до ортонормирана база на $\mathbb
R^n$. Полагаме $\bar Y = O\bar Z$. От теорията на многомерното
нормално разпределение, линейно (ортогонално) преобразувание на
вектор от независими $N(0,1)$ величини отново дава вектор от
независими $N(0,1)$ величини (приемаме този факт за известен -
той следва от инвариантността на плътността $\prod_j\frac1{\sqrt
{2\pi}}e^{-z_j^2/2}$ спрямо ротации, тъй като зависи само от нормата
$\|\bar z\|^2=\sum z_j^2$, която $O$ запазва). Следователно $Y_1,
\dots,Y_n$ са независими, $Y_i\sim N(0,1)$.

Понеже $O$ е ортогонална, тя запазва евклидовата норма:
$\sum_{j=1}^n Z_j^2 = \|\bar Z\|^2 = \|\bar Y\|^2 = \sum_{j=1}^n
Y_j^2 = Y_1^2 + \sum_{j=2}^n Y_j^2$. От избора на първия ред на $O$:
\[
  Y_1 = \frac1{\sqrt n}\sum_{j=1}^n Z_j = \frac1{\sqrt n}\sum_{j=1}^n \frac{X_j-\mu}\sigma = \frac{\sqrt n(\hat\mu-\mu)}\sigma,
  \qquad\text{т.е.}\qquad Y_1^2 = n\frac{(\hat\mu-\mu)^2}{\sigma^2}.
\]
Сравнявайки с~разлагането по-горе, $\sum_{j=2}^n Y_j^2 = \sum_j
Z_j^2 - Y_1^2 = \tfrac{(n-1)S^2}{\sigma^2}$. Понеже $Y_1$ е функция
единствено на $\bar Z$ през първия ред на $O$, а~$\sum_{j=2}^nY_j^2$
е функция единствено на $Y_2,\dots,Y_n$, и~$Y_1\perp\!\!\!\perp
(Y_2,\dots,Y_n)$ (независимост на компонентите на $\bar Y$), следва
\[
  \hat\mu = \mu+\frac{\sigma Y_1}{\sqrt n}\ \perp\!\!\!\perp\ \frac{(n-1)S^2}{\sigma^2} = \sum_{j=2}^n Y_j^2,
\]
т.е.\ $\hat\mu\perp\!\!\!\perp S^2$. Освен това $Y_2,\dots,Y_n$ са
$n-1$ на брой независими $N(0,1)$ величини, следователно
$\sum_{j=2}^n Y_j^2 = \tfrac{(n-1)S^2}{\sigma^2}\sim\chi^2(n-1)$
(отново по Тема~11). \qedbg
\end{proofbox}

\begin{thmbox}[~(доверителен интервал за \texorpdfstring{$\mu$}{мю} при неизвестно \texorpdfstring{$\sigma^2$}{сигма^2})]{Теорема}
Нека $X\sim N(\mu,\sigma^2)$, $\sigma^2$ неизвестно, и~$\hat\mu,S^2$
както по-горе. Тогава доверителният интервал за $\mu$ с~ниво на
доверие $\gamma$ е
\[
  \Big(\hat\mu-q\frac S{\sqrt n},\ \hat\mu+q\frac S{\sqrt n}\Big), \qquad q\ \text{е } \tfrac{1+\gamma}2\text{-ти квантил на } t(n-1).
\]
\end{thmbox}

\begin{proofbox}
Разглеждаме
\[
  T = \frac{\hat\mu-\mu}{S/\sqrt n} = \frac{\dfrac{\hat\mu-\mu}{\sigma/\sqrt n}}{\dfrac{S}{\sigma}} = \frac{\dfrac{\hat\mu-\mu}{\sigma/\sqrt n}}{\sqrt{\dfrac{(n-1)S^2/\sigma^2}{n-1}}}.
\]
Числителят е $\sim N(0,1)$ (по първата теорема), знаменателят е
$\sqrt{\chi^2(n-1)/(n-1)}$ (по предната теорема), и~двете са
независими (защото $\hat\mu\perp\!\!\!\perp S^2$). По дефиницията на
$t$-разпределение (Тема~11), $T\sim t(n-1)$ - централна статистика
(разпределението $t(n-1)$ не зависи от $\mu$). Избираме симетрични
квантили $\pm q$ на $t(n-1)$ с~$\mathbb P(-q\le T\le q)=\gamma$
(т.е.\ $q$ е $\tfrac{1+\gamma}2$-тият квантил на $t(n-1)$, по
симетрията на разпределението на Стюдънт). Аналогично на предната
теорема:
\[
  \gamma = \mathbb P\Big(-q\le\frac{\hat\mu-\mu}{S/\sqrt n}\le q\Big) = \mathbb P\Big(\hat\mu-q\frac S{\sqrt n}\le\mu\le\hat\mu+q\frac S{\sqrt n}\Big). \qedbg
\]
\end{proofbox}

\section{Доверителни интервали за дисперсията на нормално разпределение}

\begin{thmbox}[~(доверителен интервал за \texorpdfstring{$\sigma^2$}{сигма^2} при известно \texorpdfstring{$\mu$}{мю})]{Теорема}
Нека $X\sim N(\mu,\sigma^2)$, $\mu$ известно, $\hat\sigma^2=\tfrac1n
\sum_{j=1}^n(X_j-\mu)^2$. Тогава доверителният интервал за $\sigma^2$
с~ниво на доверие $\gamma$ е
\[
  \Big(\frac{n\hat\sigma^2}{q_2},\ \frac{n\hat\sigma^2}{q_1}\Big),
\]
където $q_1<q_2$ са квантили на $\chi^2(n)$, отрязващи по
$\tfrac{1-\gamma}2$ вероятност във всяка опашка.
\end{thmbox}

\begin{proofbox}
$T=n\hat\sigma^2/\sigma^2 = \sum_{j=1}^n\big(\tfrac{X_j-\mu}\sigma
\big)^2 \sim\chi^2(n)$ (Тема~11) е централна статистика, намаляваща
функция на $\theta=\sigma^2$. Избираме $q_1<q_2$ с~$\mathbb P(q_1\le
T\le q_2)=\gamma$. Тогава (,,обръщайки`` неравенствата спрямо
$\sigma^2$, което разменя посоката поради намаляващата зависимост):
\[
  q_1\le\frac{n\hat\sigma^2}{\sigma^2} \iff \sigma^2\le\frac{n\hat\sigma^2}{q_1},
  \qquad
  \frac{n\hat\sigma^2}{\sigma^2}\le q_2 \iff \sigma^2\ge\frac{n\hat\sigma^2}{q_2},
\]
следователно $\gamma = \mathbb P\big(\tfrac{n\hat\sigma^2}{q_2}\le
\sigma^2\le\tfrac{n\hat\sigma^2}{q_1}\big)$. \qedbg
\end{proofbox}

\begin{thmbox}[~(доверителен интервал за \texorpdfstring{$\sigma^2$}{сигма^2} при неизвестно \texorpdfstring{$\mu$}{мю})]{Теорема}
Нека $X\sim N(\mu,\sigma^2)$, $\mu$ неизвестно, $S^2$ както по-горе.
Тогава доверителният интервал за $\sigma^2$ с~ниво на доверие
$\gamma$ е
\[
  \Big(\frac{(n-1)S^2}{q_2},\ \frac{(n-1)S^2}{q_1}\Big),
\]
където $q_1<q_2$ са квантили на $\chi^2(n-1)$, отрязващи по
$\tfrac{1-\gamma}2$ вероятност във всяка опашка.
\end{thmbox}

\begin{proofbox}
Идентично на предната теорема, но с~централна статистика $T=(n-1)
S^2/\sigma^2\sim\chi^2(n-1)$ (теоремата за разлагането на
квадратичната сума). \qedbg
\end{proofbox}

\section{Проверка на хипотези}

\begin{rembox}[Схема на Нойман--Пирсън]
\begin{enumerate}[label=\arabic*)]
  \item Допускаме, че хипотезата (,,базовата``/,,нулевата``) е вярна.
  \item Конструираме събитие (\textbf{критична област}) $W\subseteq\mathbb R^n$ ($n$ - размерността на извадката); ако $\bar X\in W$, отхвърляме хипотезата, иначе я приемаме.
\end{enumerate}
Изборът на $W$ зависи от вероятността за грешка, която сме готови да
допуснем.
\end{rembox}

\begin{defbox}[~(Хипотези, грешки от I и~II род)]
Нека $X\sim F_X(x;\theta)$, $\theta\in\mathbb R^s$, и~$\bar X=(X_1,
\dots,X_n)$ са i.i.d.\ наблюдения над $X$. Разглеждаме:
\[
  H_0:\ \theta=\theta_0 \quad(\text{нулева/базова хипотеза}),
  \qquad
  H_1:\ \theta=\theta_1 \quad(\text{алтернативна хипотеза}).
\]
(По-общо, $H_0,H_1$ могат да са \emph{сложни} - напр.\ $\theta<
\theta_0$, $\theta>\theta_0$, $\theta\neq\theta_0$ - за разлика от
\emph{простите} хипотези от вида $\theta=\theta_i$.) Търсим критична
област $W\subseteq\mathbb R^n$: ако $\bar X\in W$, отхвърляме $H_0$;
иначе приемаме $H_0$. Дефинираме:
\[
  \alpha_W = \mathbb P(\bar X\in W \mid H_0) \quad(\textbf{грешка от I род}),
\]
\[
  \beta_W = \mathbb P(\bar X\notin W \mid H_1) \quad(\textbf{грешка от II род}).
\]
\end{defbox}

\begin{rembox}[Монотонност и~компромис]
Ако $W_1\subseteq W_2$, то $\alpha_{W_1}\le\alpha_{W_2}$ и~$\beta_{W_1}\ge\beta_{W_2}$ - намаляването на едната грешка (чрез
свиване на критичната област) неизбежно увеличава другата. По
конвенция, обикновено грешката от I род (погрешно отхвърляне на
вярна $H_0$) се смята за по-критична, затова фиксираме малка граница
за $\alpha$ (напр.\ $0.05$) и~се стремим да минимизираме $\beta$ при
това ограничение.
\end{rembox}

\begin{defbox}[~(Оптимална критична област)]
Нека разглеждаме проста срещу проста хипотеза и~нека за фиксирана
грешка от първи род $\alpha$ е в~сила $\alpha=\alpha_{W^*}$ за дадено
$W^*\subseteq\mathbb R^n$. Тогава $W^*$ е \textbf{оптимална критична
област} за $\alpha$, ако
\[
  \beta_{W^*} = \min_{W:\ \alpha_W=\alpha} \beta_W,
\]
т.е.\ $W^*$ минимизира вероятността за грешка от втори род измежду
всички критични области със същата грешка от първи род.
\end{defbox}

\begin{thmbox}[~(на Нойман--Пирсън)]{Лема}
Нека $H_0:\theta=\theta_0$ срещу $H_1:\theta=\theta_1$ (проста
срещу проста хипотеза), с~функции на правдоподобие $L_0(\bar x)=
L(\bar x,\theta_0)$, $L_1(\bar x)=L(\bar x,\theta_1)$. Нека
$W^*\subseteq\mathbb R^n$ е такова, че за някое $\kappa>0$:
\[
  W^* \subseteq \{\bar x\in\mathbb R^n : L_1(\bar x)\ge\kappa L_0(\bar x)\},
  \qquad
  (W^*)^c \subseteq \{\bar x\in\mathbb R^n : L_1(\bar x)\le\kappa L_0(\bar x)\}.
\]
Тогава $W^*$ е оптимална критична област за $\alpha=\alpha_{W^*}$.
\end{thmbox}

\begin{proofbox}
Нека $W\subseteq\mathbb R^n$ е произволна критична област с~$\alpha_W
=\alpha_{W^*}=\alpha$, и~нека $\beta=\beta_W$, $\beta^*=\beta_{W^*}$.
Целта е да покажем $\beta^*\le\beta$.

Разбиваме $W^c$ и~$(W^*)^c$ по двойки дизюнктни части:
\[
  W^c = \big(W^c\cap(W^*)^c\big)\cup\big(W^c\cap W^*\big),
  \qquad
  (W^*)^c = \big(W^c\cap(W^*)^c\big)\cup\big(W\cap(W^*)^c\big).
\]
Интегрирайки $L_1$ по тези разбивания:
\[
  \beta = \int_{W^c} L_1\,d\bar x = \int_{W^c\cap(W^*)^c}L_1\,d\bar x + \int_{W^c\cap W^*}L_1\,d\bar x,
\]
\[
  \beta^* = \int_{(W^*)^c}L_1\,d\bar x = \int_{W^c\cap(W^*)^c}L_1\,d\bar x + \int_{W\cap(W^*)^c}L_1\,d\bar x.
\]
Изваждайки:
\[
  \beta - \beta^* = \int_{W^c\cap W^*}L_1\,d\bar x - \int_{W\cap(W^*)^c}L_1\,d\bar x.
\]
От условията на лемата: $W^c\cap W^*\subseteq W^*\subseteq\{L_1\ge
\kappa L_0\}$, следователно $L_1\ge\kappa L_0$ навсякъде в~$W^c\cap
W^*$; и~$W\cap(W^*)^c\subseteq(W^*)^c\subseteq\{L_1\le\kappa L_0\}$,
следователно $L_1\le\kappa L_0$ навсякъде в~$W\cap(W^*)^c$. Тогава
\[
  \beta-\beta^* \ge \int_{W^c\cap W^*}\kappa L_0\,d\bar x - \int_{W\cap(W^*)^c}\kappa L_0\,d\bar x = \kappa\left(\int_{W^c\cap W^*}L_0\,d\bar x - \int_{W\cap(W^*)^c}L_0\,d\bar x\right).
\]
Остава да покажем, че скобата е нула. Понеже $\alpha_W=\alpha_{W^*}$:
\[
  \int_W L_0\,d\bar x = \alpha_W = \alpha_{W^*} = \int_{W^*}L_0\,d\bar x.
\]
Разбиваме отново по дизюнктни части: $W=(W\cap W^*)\cup(W\cap
(W^*)^c)$ и~$W^*=(W\cap W^*)\cup(W^*\cap W^c)$, откъдето
\begin{align*}
  \int_{W\cap W^*}L_0 + \int_{W\cap(W^*)^c}L_0 &= \int_W L_0 = \int_{W^*}L_0 = \int_{W\cap W^*}L_0 + \int_{W^*\cap W^c}L_0\\
  &\quad\Longrightarrow\quad \int_{W\cap(W^*)^c}L_0 = \int_{W^*\cap W^c}L_0,
\end{align*}
т.е.\ точно скобата по-горе е нула (защото $W^*\cap W^c=W^c\cap
W^*$). Следователно $\beta-\beta^*\ge\kappa\cdot0=0$, т.е.\
$\beta\ge\beta^*$ - $W^*$ минимизира грешката от втори род измежду
критичните области със същата грешка от първи род, следователно е
оптимална. \qedbg
\end{proofbox}
